\chapter{Ordered Sets, Treaps and Other Structures}\label{ch:advds}

Segment trees answer questions about a fixed array. This chapter collects the structures for the other situations: a set whose
contents change and whose \emph{order} matters (the neighbours of a value, the $k$-th smallest), a sequence whose \emph{shape} changes
(insertions in the middle, reversals, cut and paste), many versions of a sequence at once, offline questions in three dimensions, and
range questions about \emph{values} rather than positions. It ends with hashing, the right tool when no order is needed at all, and
with a way to hash whole sets.

\begin{chapterintro}{\gold\ \plat\ \ioi}{Chapter~\ref{ch:cpp} (\cc{set}, \cc{map} and the policy-based tree),
Chapter~\ref{ch:ds-basics} (interval sets and linked lists) and Chapter~\ref{ch:segtree} (Fenwick trees and persistence).}
  \item combine sorted containers to answer neighbour, gap and order-statistic questions in $\Oh{\log n}$;
  \item use treaps keyed by position to insert, delete, reverse and move whole ranges, and prove their expected depth;
  \item keep every version of a sequence with a persistent treap, and budget its memory;
  \item count three-dimensional dominance offline with CDQ divide and conquer;
  \item answer $k$-th smallest and counting queries on ranges with a wavelet matrix;
  \item choose between ordered containers, hash tables, sorting and plain arrays from measured costs, and hash whole sets.
\end{chapterintro}

% ======================================================================================================================
\section[Sets beyond lookups]{Sets beyond lookups \hfill \gold}\label{sec:setsorder}

A \cc{set}, \cc{multiset} or \cc{map} keeps its elements sorted in a balanced binary search tree, so besides insert, erase and
find it answers \emph{ordered} questions in $\Oh{\log n}$: Chapter~\ref{ch:cpp} showed how \cc{lower_bound} walks the tree
(Figure~\ref{fig:bst-lower}), and Section~\ref{sec:ordered-sets} kept a set of disjoint intervals. Table~\ref{tab:setops} collects
the idioms that come up again and again. Two habits prevent most bugs: compare an iterator with \cc{end()} (or \cc{begin()}) before
dereferencing it or its predecessor, and erase \emph{one} copy from a multiset with \cc{ms.erase(ms.find(x))}, because
\cc{ms.erase(x)} removes all of them.

The flexibility has a price in memory and in locality. Every element of a set is a separate allocation that holds the key, three
pointers and a colour: on the machine used for this book, a \cc{set<int>} of $10^6$ random numbers occupied 51.5\,MB, thirteen times the
4\,MB of the numbers themselves, and $10^6$ insertions followed by $10^6$ lookups took 1.14\,s. Sorting a \cc{vector} and searching it
with \cc{binary_search} did the same work in 0.14\,s and 5.3\,MB. When the collection does not change between queries, a sorted vector
is the right structure; a set earns its cost only when insertions and queries interleave.

\begin{table}[htbp]
\centering\small
\begin{tabularx}{\linewidth}{@{}>{\ttfamily\raggedright\arraybackslash}p{5.6cm}X@{}}
\toprule
\normalfont expression & \normalfont meaning \\
\midrule
*s.lower\_bound(x) & the smallest element $\ge x$ (check \texttt{!= s.end()} first) \\
*s.upper\_bound(x) & the smallest element $> x$ \\
*prev(s.lower\_bound(x)) & the largest element $< x$ (check \texttt{!= s.begin()} first) \\
*prev(s.upper\_bound(x)) & the largest element $\le x$ \\
*s.begin(), *s.rbegin() & the minimum and the maximum \\
next(it), prev(it) & the neighbours of an element, amortised $\Oh{1}$ \\
ms.erase(ms.find(x)) & remove \emph{one} copy of $x$ from a multiset \\
s.emplace\_hint(s.end(), x) & append a value larger than all others, amortised $\Oh{1}$ \\
s.extract(it) & take the node out without freeing it (C++17); change its key through \texttt{.value()} (\texttt{.key()} in a
  map) and \texttt{insert} it back, with no new allocation \\
\bottomrule
\end{tabularx}
\caption{Idioms for \cc{std::set}, \cc{multiset} and \cc{map}; each costs $\Oh{\log n}$ unless stated otherwise.}
\label{tab:setops}
\end{table}

The interesting problems keep \emph{two} containers that describe the same situation from two sides, and update both on every event.

\begin{problem}{Traffic Lights (CSES)}
A street of length $x$ receives traffic lights one at a time, at distinct positions; after each, report the length of the longest
stretch of the street without lights ($x \le 10^9$, $n \le 2 \cdot 10^5$).
\solution Keep the positions in a \cc{set}, with the two ends of the street as sentinels, and the lengths of the stretches between
neighbouring lights in a \cc{multiset}. A new light at $p$ lies inside exactly one stretch, between the first light after $p$
(\cc{upper_bound}) and its predecessor. That stretch is replaced by two (Figure~\ref{fig:traffic}), and the answer is the largest
element of the multiset. Each light costs a constant number of $\Oh{\log n}$ operations.
\end{problem}

\begin{figure}[htbp]
\centering
\begin{tikzpicture}[x=1.1cm]
  \foreach \row/\lights/\gaps/\ttx in {0/{0,3,8}/{3/1.5, 5/5.5}/{light at 3: longest 5},
      1/{0,3,6,8}/{3/1.5, 3/4.5, 2/7}/{light at 6: longest 3},
      2/{0,2,3,6,8}/{2/1, 1/2.5, 3/4.5, 2/7}/{light at 2: longest 3}} {
    \begin{scope}[yshift=-1.05*\row cm]
      \draw[line width=1pt, draw=fgline] (0, 0) -- (8, 0);
      \foreach \pp in \lights \fill[fgverm] (\pp, 0) circle (2.6pt);
      \foreach \glen/\gpos in \gaps \node[font=\sffamily\scriptsize\bfseries, text=fgblue, above=2pt] at (\gpos, 0) {\glen};
      \node[lbl, anchor=west] at (8.4, 0) {\ttx};
    \end{scope}
  }
  \draw[line width=0.5pt, fgink!60] (0, -2.45) -- (8, -2.45);
  \foreach \pp in {0, ..., 8} {
    \draw[line width=0.5pt, fgink!60] (\pp, -2.38) -- (\pp, -2.52);
    \node[lbl, below=1pt] at (\pp, -2.5) {\pp};
  }
\end{tikzpicture}
\caption{Traffic Lights on a street of length 8 with lights added at 3, 6 and 2. The set holds the red points, the multiset the
blue lengths; each new light removes one length and inserts two.}
\label{fig:traffic}
\end{figure}

\begin{code}
void solve() {
    ll x;
    int n;
    cin >> x >> n;
    set<ll> lights{0, x};  // the two ends of the street act as lights
    multiset<ll> gaps{x};  // the lengths between neighbouring lights
    vi ans;
    rep(i, n) {
        ll p;
        cin >> p;
        auto it = lights.upper_bound(p);  // the light just right of p
        ll right = *it, left = *prev(it);
        gaps.erase(gaps.find(right - left));  // one gap splits into two
        gaps.insert(p - left), gaps.insert(right - p);
        lights.insert(p);
        ans.push_back(*gaps.rbegin());
    }
    print(ans);
}
\end{code}

The same problem has an offline solution without any tree. Run time backwards: start from all the lights, sorted once, and
\emph{remove} them in reverse order. A removal merges two stretches into one, so the longest stretch can only grow, and a running
maximum together with a linked list over the sorted positions (Section~\ref{sec:linked}) answers every step in $\Oh{1}$. Reversing
time to turn splits into merges is a recurring trick. The two-container idea also recurs: Section~\ref{sec:sliding} keeps the lower
and the upper half of a sliding window in two multisets to maintain its median.

\paragraph{Custom comparators.} A set orders its elements by any strict weak ordering. Two elements of which neither is ``less'' than
the other count as \emph{equal}, and only one of them is kept, even if they differ in fields that the comparator ignores. A set of
segments ordered by right end alone would silently drop every segment whose right end is already present; a unique id as the last
tie-breaker keeps them apart.

\begin{code}
struct Seg {
    int l, r, id;
};
struct ByRight {  // order by right end; the id keeps distinct segments apart
    bool operator()(const Seg& a, const Seg& b) const {
        return make_pair(a.r, a.id) < make_pair(b.r, b.id);
    }
};
set<Seg, ByRight> active;
\end{code}

\begin{pitfall}[Comparators that change]
Never modify a field that the comparator uses while the element is inside the set: the tree does not notice, and later searches go
the wrong way. Take the node out with \cc{extract}, change it, and insert it again. For the same reason a comparator must not depend
on state that changes, such as a global ``current $x$'', unless you can prove that the relative order of the stored elements never
changes; the sweep-line algorithms that order segments by their height at the current $x$ (Chapter~\ref{ch:geometry}) rest on
exactly such a proof.
\end{pitfall}

% ======================================================================================================================
\section[Order statistics]{Order statistics \hfill \gold}\label{sec:orderstats}

``What is the $k$-th smallest element?'' and ``how many elements are smaller than $x$?'' in a changing set are beyond
\cc{std::set}, whose iterators move one step at a time. The policy-based tree of Section~\ref{sec:pbds} stores subtree sizes:
\cc{find_by_order(k)} returns an iterator to the $k$-th smallest element (0-based) and \cc{order_of_key(x)} counts the elements
$< x$, both in $\Oh{\log n}$. When all values that will ever be inserted are known in advance, a Fenwick tree over their compressed
ranks does the same with a smaller constant: \cc{order_of_key} is a prefix sum and \cc{find_by_order} a descent
(Figure~\ref{fig:fendescent}).

\begin{problem}{Josephus Problem II (CSES)}
Children $1, \dots, n$ stand in a circle; repeatedly, $k$ children are skipped and the next one leaves the circle. Print the order
in which they leave ($n \le 2 \cdot 10^5$, $0 \le k \le 10^9$).
\solution The linked list that solved the case $k = 1$ (Figure~\ref{fig:josephus}) would walk $k$ steps per removal. Instead keep
the remaining children in an ordered set, and the index, in the current circle, of the child where counting starts. Skipping $k$
children moves the index by $k$ modulo the number of children left; \cc{find_by_order} finds the child to remove, and after it is
erased the next child has the same index (Figure~\ref{fig:josephus2}). The simulation costs $\Oh{n \log n}$ and runs in 0.16\,s for
$n = 2 \cdot 10^5$.
\end{problem}

\begin{figure}[htbp]
\centering
\begin{tikzpicture}
  \foreach \cc/\ord [count=\ii from 0] in {1/6, 2/3, 3/1, 4/7, 5/5, 6/2, 7/4} {
    \pgfmathsetmacro{\ang}{90 - 360/7*\ii}
    \node[vtx, minimum size=7mm] (J\cc) at (\ang:1.6) {\cc};
    \node[font=\sffamily\scriptsize\bfseries, text=fgverm] at (\ang:2.3) {\ord};
  }
  \node[notel, anchor=west, text width=7.4cm] at (3.0, 0) {$n = 7$, $k = 2$. The red numbers say when each child leaves:
    the order is 3, 6, 2, 7, 5, 1, 4. The indices used are $(0 + 2) \bmod 7 = 2$ (child 3), then $(2 + 2) \bmod 6 = 4$ (child 6
    among 1, 2, 4, 5, 6, 7), then $(4 + 2) \bmod 5 = 1$ (child 2), and so on. Each removal is one $\Oh{\log n}$ search, however
    large $k$ is.};
\end{tikzpicture}
\caption{Josephus Problem II with $n = 7$ and $k = 2$: the index of the next child to leave, modulo the number of children left.}
\label{fig:josephus2}
\end{figure}

\begin{code}
void solve() {  // ordered_set (cposet) is pasted above solve()
    int n, k;
    cin >> n >> k;
    ordered_set<int> s;
    rep(i, 1, n + 1) s.insert(i);
    vi out;
    int pos = 0;  // the index, in the current circle, to start counting from
    while (!s.empty()) {
        pos = (pos + k) % sz(s);  // skip k children
        auto it = s.find_by_order(pos);
        out.push_back(*it);
        s.erase(it);  // the next child now has the same index
    }
    print(out);
}
\end{code}

\usetemplate{Ordered set (PBDS)}{Ctrl+Alt+T, Shift+O}{cposet}

\paragraph{Multisets.} The tree is a set, so a value inserted twice is stored once. Store pairs \cc|{value, id}| with a unique id
(the time of insertion, for example); then \cc|order_of_key({x, -1})| counts the elements $< x$, and \cc|order_of_key({x, INF})|
those $\le x$.

\begin{pitfall}[The \texttt{less\_equal} trick]
A trick seen in many published solutions declares the tree with \cc{less_equal<T>} so that it keeps duplicates. That comparator is
not a strict ordering, and the tree then misbehaves: \cc{find} never finds anything, so \cc{s.erase(x)} does nothing, and
\cc{lower_bound} and \cc{upper_bound} swap their meanings. Use pairs.
\end{pitfall}

% ======================================================================================================================
\section[Treaps]{Treaps \hfill \plat}\label{sec:treap}

A \defn{treap} is a binary search tree in which every node also carries a random \emph{priority}, and which is also a heap by
priority: every node's priority exceeds its children's. Draw each node at the point (position in sorted order, priority), as in
Figure~\ref{fig:treap}. The search-tree property says that a node's left subtree lies to its left and its right subtree to its
right; the heap property says that every node lies above its descendants. Given the keys and the priorities, these two rules
determine the tree completely: the root must be the highest point, the left subtree is built from the points to its left, and the
right subtree from the points to its right. Everything is done with two operations:
\begin{itemize}
  \item \textbf{split}$(t, k)$ cuts $t$ into a treap of its first $k$ elements and a treap of the rest. It walks down from the root
        along the cut: a node whose left subtree holds at least $k$ elements goes to the right part, together with its right
        subtree, and the walk continues to the left; otherwise the node goes to the left part with its left subtree, and the walk
        continues to the right. Only the nodes on this path are relinked.
  \item \textbf{merge}$(a, b)$ joins two treaps where every element of $a$ comes before every element of $b$. The higher of the two
        roots becomes the root; if it is $a$'s root, its left subtree stays and its right subtree is merged with $b$, and
        symmetrically.
\end{itemize}

\begin{figure}[htbp]
\centering
\begin{tikzpicture}[x=0.8cm, y=0.4cm]
  \tikzset{tn/.style={vtx, minimum size=6.2mm, font=\sffamily\small}}
  \foreach \sx/\ttx in {0/{the treap}, 9/{split after five elements}} {
    \draw[->, line width=0.6pt, fgink!70] (\sx - 0.7, 0) -- (\sx - 0.7, 10.8) node[above, lbl] {priority};
    \foreach \yv in {0, 3, 6, 9} {
      \draw[line width=0.5pt, fgink!70] (\sx - 0.77, \yv) -- (\sx - 0.63, \yv);
      \node[lbl, anchor=east] at (\sx - 0.8, \yv) {\yv};
    }
    \draw[line width=0.6pt, fgink!70] (\sx - 0.7, 0) -- (\sx + 6.5, 0);
    \foreach \px in {0, ..., 6} \node[lbl] at (\sx + \px, -1.0) {\px};
    \node[lbl] at (\sx + 3, -2.3) {position in the sequence};
    \node[title] at (\sx + 3, 12.6) {\ttx};
    \draw[line width=1.3pt, draw=fgverm, dashed] (\sx + 4.5, -0.5) -- (\sx + 4.5, 10.6);
    \foreach \nm/\px/\pr in {a/0/3, b/1/7, c/2/2, d/3/9, e/4/4, f/5/6, g/6/1} {
      \draw[line width=0.5pt, draw=fgline, densely dotted] (\sx + \px, 0) -- (\sx + \px, \pr);
      \node[tn] (N\sx\nm) at (\sx + \px, \pr) {\nm};
    }
  }
  \draw[ed] (N0d) -- (N0b) (N0b) -- (N0a) (N0b) -- (N0c) (N0f) -- (N0g);
  \draw[hed] (N0d) -- (N0f) (N0f) -- (N0e);
  \draw[ed] (N9d) -- (N9b) (N9b) -- (N9a) (N9b) -- (N9c) (N9f) -- (N9g);
  \draw[ged] (N9d) -- (N9e);
\end{tikzpicture}
\caption{A treap holding the sequence a, \dots, g, drawn at (position, priority): parents lie above their children, and the
in-order walk reads the sequence from left to right. Splitting after five elements walks down the path d, f, e across the cut
(the red edges) and relinks it: e becomes the right child of d (green), and f, with g, forms the second treap. Merging the two
treaps gives back the first picture.}
\label{fig:treap}
\end{figure}

The template stores the nodes in a vector (node 0 is the empty tree) and keeps in every node the size of its subtree, which is what
lets split count elements. Both operations are short recursions:

\begin{code}
pair<int, int> split(int x, int k) {  // first k elements | the rest
    if (!x) return {0, 0};
    push(x);  // apply a pending lazy flag before looking at the children
    if (t[t[x].l].size >= k) {
        auto [a, b] = split(t[x].l, k);
        t[x].l = b, pull(x);  // pull recomputes size and sum from the children
        return {a, x};
    }
    auto [a, b] = split(t[x].r, k - t[t[x].l].size - 1);
    t[x].r = a, pull(x);
    return {x, b};
}
int merge(int a, int b) {  // every element of a comes before every element of b
    if (!a || !b) return a ? a : b;
    if (t[a].pri > t[b].pri) {  // the higher root stays on top
        push(a);
        t[a].r = merge(t[a].r, b), pull(a);
        return a;
    }
    push(b);
    t[b].l = merge(a, t[b].l), pull(b);
    return b;
}
\end{code}

Both functions spend $\Oh{1}$ per level of the treap that they walk, so their cost is bounded by its depth, and with random
priorities the expected depth is logarithmic.

\begin{fact}[Expected depth of a treap]
Let $i$ and $j$ be two nodes, and consider the highest node among those from $i$ to $j$ in sorted order. The searches for $i$ and for
$j$ from the root take the same turns until they reach that node, because both lie on the same side of every higher node. So $j$
is an ancestor of $i$ exactly when $j$ itself is the highest node of the range; if a node strictly between them is higher, it
separates them (in Figure~\ref{fig:treap}, d separates c from e). With independent random priorities, every node of the range is
equally likely to be the highest, so the probability is $1/(|i - j| + 1)$, and the expected number of ancestors of $i$ is
\[
  \sum_{j < i} \frac{1}{i - j + 1} + \sum_{j > i} \frac{1}{j - i + 1} = (H_i - 1) + (H_{n - i + 1} - 1) \le 2 \ln n,
\]
where $H_m = 1 + \frac12 + \dots + \frac1m \le 1 + \ln m$. Measured with the template for $n = 2 \cdot 10^5$: the average depth is
about 22, as the formula predicts, and the deepest node lies at depth 42 to 44, depending on the run.
\end{fact}

\begin{figure}[htbp]
\centering
\begin{tikzpicture}
  \node[blbl, anchor=east] at (-0.2, 0) {sequence};
  \arrs{S}{0}{0}{a/n, b/n, c/n, d/cool, e/cool, f/cool, g/n}
  \idx{S}{7}
  \draw[line width=1.6pt, draw=fgverm] ($(S-3.north east)+(0, 0.2)$) -- ($(S-3.south east)+(0, -0.2)$);
  \draw[line width=1.6pt, draw=fgverm] ($(S-6.north east)+(0, 0.2)$) -- ($(S-6.south east)+(0, -0.2)$);
  \node[blbl, anchor=east] at (-0.2, -1.5) {after};
  \arrs{T}{0}{-1.5}{a/n, b/n, c/n, f/cool, e/cool, d/cool, g/n}
  \node[notel, anchor=west, text width=8.0cm] at (5.5, -0.75) {Reversing positions $[3, 6)$: split off the first 3 elements,
    split the next 3 from the rest, flip the \emph{reversed} flag of the middle piece's root, and merge the three pieces back.
    Four $\Oh{\log n}$ operations, however long the range is.};
\end{tikzpicture}
\caption{A range reversal on an implicit treap: two splits, one flag, two merges.}
\label{fig:treaprev}
\end{figure}

\paragraph{Implicit keys.} If no key is stored at all and a node's position is defined as \emph{the number of nodes before it} in the
in-order walk, the treap represents a sequence, and split by count cuts it into a prefix and a suffix. Then every range operation is
``split off the range, act on its root, merge back'' (Figure~\ref{fig:treaprev}): a range sum is the middle piece's aggregate, a
cut-and-paste merges the three pieces in another order, and reversing a range sets a lazy \emph{reversed} flag on the middle piece's
root. The flag is pushed down only when a later split or merge walks into that node: \cc{push} swaps the node's children and passes
the flag on to them. Since the in-order walk visits the left subtree first, swapping the children of every node of a subtree reverses
the subtree's sequence.

\begin{problem}{Reversals and Sums (CSES)}
An array of $n \le 2 \cdot 10^5$ numbers receives $m \le 2 \cdot 10^5$ operations: reverse a subarray, or report the sum of a
subarray.
\solution An implicit treap whose nodes store the sum of their subtree. The template's \cc{reverse_range} and \cc{range_sum} are the
two splits and two merges of Figure~\ref{fig:treaprev}.
\end{problem}

\begin{code}
void solve() {  // Treap (cptreap) is pasted above solve()
    int n, m;
    cin >> n >> m;
    Treap tr;
    rep(i, n) {
        ll x;
        cin >> x;
        tr.insert(i, x);
    }
    while (m--) {
        int type, a, b;
        cin >> type >> a >> b;
        if (type == 1) tr.reverse_range(a - 1, b);  // split, flag, merge
        else print(tr.range_sum(a - 1, b));
    }
}
\end{code}

\usetemplate{Implicit treap}{Ctrl+Alt+Y, T}{cptreap}

With $n = m = 2 \cdot 10^5$ the program takes 0.56\,s. A treap pays for its flexibility with a constant several times larger than a
segment tree's: every operation walks about 22 levels of nodes scattered through memory, and does so four times. Substring Reversals
and Cut and Paste are the same skeleton with characters, and a new kind of operation needs only a new aggregate in \cc{pull} and,
for lazy changes, a new flag in \cc{push}.

\begin{recognize}
The array itself changes shape: ``move the segment $[l, r)$ to the front'', ``reverse a range'', ``insert an element at position
$p$'', rope-like string editing. If only the values change, a segment tree is simpler and faster; if elements are only inserted and
all operations are known in advance, processing them backwards with a Fenwick descent often avoids the treap altogether.
\end{recognize}

\subsection{Persistent treaps}\label{sec:ptreap}

Split and merge change only the nodes on the paths they walk, so they can \emph{copy} those nodes instead of modifying them, exactly
like path copying in a segment tree (Section~\ref{sec:persistent}). Every operation then returns a new root and leaves all old
versions intact: a persistent sequence. Concatenating two versions shares all their nodes, so a copy and paste of a whole block costs
$\Oh{\log n}$ new nodes, even when the text doubles in length with every paste.

Stored priorities do not survive this. When a version is pasted into itself, the same nodes, with the same priorities, appear twice,
so priorities are no longer independent, the depth analysis above no longer applies, and long chains of equal priorities can form.
The template therefore stores no priorities at all: \cc{merge}$(a, b)$ puts $a$'s root on top with probability $|a| / (|a| + |b|)$.
That is exactly the probability that the highest of $|a| + |b|$ independent priorities lies in $a$, so every merge makes the choice
that a fresh random treap would make, whatever the history of the nodes.

\begin{problem}{Versions of a list}
Starting from an array (version 0), each operation creates a new version from an existing one, by inserting an element at a position
or deleting one, or asks for the sum of a range of an existing version.
\end{problem}

\begin{code}
void solve() {  // PTreap (cpptreap) is pasted above solve()
    int n, q;
    cin >> n >> q;
    vi a(n);
    for (ll& x : a) cin >> x;
    PTreap t;
    vi ver{t.build(a)};  // ver[k]: the root of version k; old roots stay valid
    while (q--) {
        int type, k;
        cin >> type >> k;
        if (type == 1) {  // version k with x inserted before position p
            int p;
            ll x;
            cin >> p >> x;
            ver.push_back(t.insert(ver[k], p, x));
        } else if (type == 2) {  // version k without the element at p
            int p;
            cin >> p;
            ver.push_back(t.erase(ver[k], p));
        } else {  // the sum of positions [l, r) of version k
            int l, r;
            cin >> l >> r;
            print(t.sum(ver[k], l, r));
        }
    }
}
\end{code}

\usetemplate{Persistent treap}{Ctrl+Alt+U, R}{cpptreap}

\begin{pitfall}[Memory]
Nothing is ever freed. Measured on $2 \cdot 10^5$ random operations over a list of $2 \cdot 10^5$ elements, every update created 46
new nodes of 32 bytes: 6.3 million nodes, 200\,MB. A vector that grows by doubling briefly holds both its old and its new buffer, so
the peak was near 400\,MB. Reserve the expected size up front (\cc{t.t.reserve(n + 50 * q)}), and if old versions become
unreachable, rebuild the live versions into a fresh pool every so often.
\end{pitfall}

% ======================================================================================================================
\section[CDQ divide and conquer]{CDQ divide and conquer \hfill \ioi}\label{sec:cdq}

Offline problems about points in three dimensions, such as ``for each point, count the points that are $\le$ it in all three
coordinates'', have a standard reduction named after Chen Danqi. Sort the points by the first coordinate $a$, so that a point can only
be dominated by points before it, and split the sorted list in half. A pair of points within one half is handled recursively. A pair
with one point in each half satisfies the first coordinate automatically, left before right, so for these \emph{cross pairs} only the
second and third coordinates remain: a two-dimensional dominance count, which is a sweep. Visit both halves in order of $b$; insert
each left point's $c$ into a Fenwick tree, and for each right point count the inserted values $\le c$ (Figure~\ref{fig:cdq}). If the
recursion returns each half sorted by $b$, the sweep is a merge, so each level of the recursion costs $\Oh{n \log n}$ and the whole
algorithm $\Oh{n \log^2 n}$.

\begin{figure}[htbp]
\centering
\begin{tikzpicture}[x=0.9cm, y=0.9cm]
  \fill[fgsky!22] (0, 0) rectangle (4.0, 3.0);
  \draw[line width=0.8pt, draw=fgblue, dashed] (4.0, 0) -- (4.0, 3.0) -- (0, 3.0);
  \draw[->, line width=0.6pt, fgink!70] (0, 0) -- (6.4, 0) node[right, lbl] {$b$};
  \draw[->, line width=0.6pt, fgink!70] (0, 0) -- (0, 4.4) node[above, lbl] {$c$};
  \foreach \px/\py in {1/1, 2.5/2.2, 3.4/0.6, 5.2/1.4, 1.6/3.6}
    \node[vtxs, cool] at (\px, \py) {};
  \foreach \px/\py in {0.8/3.9, 5.6/2.4}
    \node[vtxs, warn] at (\px, \py) {};
  \node[vtxs, hot] (Q) at (4.0, 3.0) {};
  \node[lbl, anchor=south west] at (Q.north east) {$(b_y, c_y)$};
  \legendat{(7.0, 3.6)}{cool/left half (smaller $a$), warn/right half, hot/the right point $y$}
  \node[notel, anchor=north west, text width=6.4cm] at (6.85, 2.3) {The cross pairs of $y$ are the left points in the shaded
    quadrant: here 3. Sweeping $b$ upwards, with the left points' $c$ values in a Fenwick tree, counts them for every right point
    in one pass.};
\end{tikzpicture}
\caption{The merge step of CDQ: between the two halves only two coordinates are left.}
\label{fig:cdq}
\end{figure}

\begin{pseudo}{CDQ divide and conquer for dominance counting}
\begin{algorithmic}[1]
\State sort the points by $(a, b, c)$ and merge identical points into one item with a multiplicity
\Procedure{CDQ}{$l, r$} \Comment{items $l, \dots, r - 1$; on return they are sorted by $(b, c)$}
  \If{$r - l \le 1$} \State \Return \EndIf
  \State $m \gets \lfloor (l + r) / 2 \rfloor$; \ \Call{CDQ}{$l, m$}; \ \Call{CDQ}{$m, r$}
  \ForAll{items $y$ of the right half, in order of $(b, c)$}
    \State add to the Fenwick tree the $c$ of the next left items with $(b, c) \le (b_y, c_y)$
    \State $\mathit{ans}[y] \gets \mathit{ans}[y] + {}$the number of added values $\le c_y$
  \EndFor
  \State remove the added values again; merge the two halves by $(b, c)$
\EndProcedure
\end{algorithmic}
\end{pseudo}

\begin{problem}{Three-dimensional dominance}
For each of $n \le 10^5$ points $(a, b, c)$, count the other points $(a', b', c')$ with $a' \le a$, $b' \le b$ and $c' \le c$.
\solution CDQ as above. Identical points dominate each other, which a rule of ``left half before right half'' would count in only
one direction, so the template first groups identical points into one item with a multiplicity: every copy is dominated by the other
copies of its group plus everything that dominates the group. The program takes 0.32\,s for $10^5$ points.
\end{problem}

\begin{code}
void solve() {  // cdq_dominance (cpcdq) is pasted above solve()
    int n;
    cin >> n;
    vector<array<ll, 3>> pts(n);
    for (auto& [a, b, c] : pts) cin >> a >> b >> c;
    for (ll d : cdq_dominance(pts)) print(d);
}
\end{code}

\usetemplate{CDQ divide and conquer (3D dominance)}{Ctrl+Alt+U, 3}{cpcdq}

\begin{trick}[CDQ beyond counting]
The ``first coordinate'' can be \emph{time}. With updates and queries interleaved, CDQ on the timeline lets every update of the left
half act on the queries of the right half through a static sweep, which turns a dynamic problem into an offline static one: points
added and rectangle sums asked in any order, for example. CDQ also speeds up DPs whose transition is a dominance condition, such as
the longest chain of points increasing in three coordinates: recurse on the left half first, let its finished values update the right
half, and only then recurse on the right half.
\end{trick}

% ======================================================================================================================
\section[Wavelet trees]{Wavelet trees \hfill \ioi}\label{sec:wavelet}

A wavelet tree answers questions about the \emph{values} in a range of positions, such as the $k$-th smallest element of $a[l, r)$ or
the number of elements of $a[l, r)$ that are $\le x$, online and in $\Oh{\log n}$ each. The compact form used in contests, the
\defn{wavelet matrix}, works on bits. Replace every value by its rank among the distinct values, a number of $\lceil \log_2 n \rceil$
bits. The top level records, for every position, the highest bit of its rank, and then stably moves the elements whose bit is 0 in
front of those whose bit is 1; the next level records the next bit in this new order and partitions again, down to the lowest bit
(Figure~\ref{fig:wavelet}). Each level is a bit vector with prefix counts of ones, so ``how many ones are there among the first $i$
elements of level $h$?'' takes $\Oh{1}$.

To find the $k$-th smallest (0-based) of a range, look at the current bit of the elements of the range. If at least $k + 1$ of them
have bit 0, the answer's bit is 0 and the range continues among the zeros of the next level; otherwise the answer's bit is 1, $k$
decreases by the number of zeros, and the range continues among the ones. The stable partition keeps the elements of the range
contiguous at every level, and both ends of the new range follow from two prefix counts, so a query costs $\Oh{1}$ per level.
Counting the elements $\le x$ walks down the bits of $x$ in the same way, adding up the zeros that it passes where $x$ has a 1.

\begin{figure}[htbp]
\centering
\begin{tikzpicture}
  \foreach \yy/\ttx in {0/{bit 2}, -1.25/{bit 1}, -2.5/{bit 0}} \node[blbl, anchor=east] at (-0.2, \yy) {\ttx};
  \arrs{W}{0}{0}{5/cool, 2/n, 6/cool, 1/n, 7/cool, 3/n, 0/n, 4/cool}
  \arrs{X}{0}{-1.25}{2/cool, 1/n, 3/cool, 0/n, 5/n, 6/cool, 7/cool, 4/n}
  \arrs{Y}{0}{-2.5}{1/cool, 0/n, 5/cool, 4/n, 2/n, 3/cool, 6/n, 7/cool}
  \foreach \ii in {1, ..., 8} {
    \pgfmathtruncatemacro{\jj}{\ii - 1}
    \node[lbl] at ($(W-\ii.north)+(0, 0.24)$) {\jj};
  }
  \cellframe{W-2}{W-6}{fgverm}
  \cellframe{X-1}{X-3}{fgverm}
  \cellframe{Y-5}{Y-6}{fgverm}
  \node[notel, anchor=west, text width=8.0cm] at (6.0, -1.25) {Ranks $5, 2, 6, 1, 7, 3, 0, 4$ at positions 0 to 7. Blue cells
    have a 1 in the bit of their level, and each level puts the elements with a 0 first. The 3rd smallest ($k = 2$) of positions
    $[1, 6)$: at bit 2 the range has three 0s and $k < 3$, so the bit is 0 and the range becomes $[0, 3)$; at bit 1 it has one 0,
    so the bit is 1, $k$ becomes 1 and the range $[4, 6)$; at bit 0 there is one 0 again, so the bit is 1. The answer is rank
    $011_2 = 3$.};
\end{tikzpicture}
\caption{A wavelet matrix over eight ranks: one bit per level, stably partitioned, with the range of a query followed down.}
\label{fig:wavelet}
\end{figure}

\begin{problem}{$k$-th smallest in a range}
Answer $q$ online queries ``the $k$-th smallest value among $a_l, \dots, a_r$'' for an array of $n$ numbers ($n, q \le 2 \cdot 10^5$,
values up to $10^9$).
\end{problem}

\begin{code}
void solve() {  // Wavelet (cpwavelet) is pasted above solve()
    int n, q;
    cin >> n >> q;
    vi a(n);
    for (ll& x : a) cin >> x;
    Wavelet wt(a);  // compresses the values: log2(n) levels, not log2(1e9)
    while (q--) {
        int l, r, k;
        cin >> l >> r >> k;
        print(wt.kth(l - 1, r, k - 1));  // the k-th smallest of a[l..r]
    }
}
\end{code}

\usetemplate{Wavelet tree}{Ctrl+Alt+Y, W}{cpwavelet}

The template packs each level into 64-bit words with a prefix count per word, so a rank query is one table lookup plus one
\cc{__builtin_popcountll}, and the 18 levels for $n = 2 \cdot 10^5$ take about 0.7\,MB, small enough to stay in the cache. The program
answers $2 \cdot 10^5$ queries in 0.28\,s; an earlier version of the template that stored a full integer count for every element and
level, 14\,MB, took 0.72\,s on the same input. Table~\ref{tab:kthrange} compares the structures that answer these queries.

\begin{table}[htbp]
\centering\small
\begin{tabular}{@{}lccl@{}}
\toprule
structure & $k$-th smallest & count $\le x$ & memory for $n = 2 \cdot 10^5$ \\
\midrule
merge sort tree (sorted lists in segment tree nodes) & $\Oh{\log^3 n}$ & $\Oh{\log^2 n}$ & 3.6 million numbers \\
persistent segment tree over the values & $\Oh{\log n}$ & $\Oh{\log n}$ & 3.8 million nodes, 61\,MB \\
wavelet matrix & $\Oh{\log n}$ & $\Oh{\log n}$ & 3.6 million bits, 0.7\,MB \\
offline: sorted queries, sweep and Fenwick tree & --- & $\Oh{\log n}$ & $\Oh{n}$ \\
\bottomrule
\end{tabular}
\caption{Online range questions about values, none of which allow updates. The persistent segment tree is
Section~\ref{sec:persistent}; when all queries are known in advance, counting needs only a sweep (Section~\ref{sec:sweepqueries}).}
\label{tab:kthrange}
\end{table}

% ======================================================================================================================
\section[Hash maps and hashing sets]{Hash maps and hashing sets \hfill \gold}\label{sec:hashmaps-adv}

\cc{unordered_map} offers expected $\Oh{1}$ operations, but as Chapter~\ref{ch:cpp} showed (Figure~\ref{fig:hashattack}), its default
hash for integers is the identity, and a test built against it makes every key collide. The seeded \cc{custom_hash} of the library
fixes that. How much faster than a \cc{map} is a hash table, and is \cc{unordered_map} the fastest one? Table~\ref{tab:dicts} gives
measured answers.

\begin{table}[htbp]
\centering\small
\begin{tabular}{@{}lrr@{}}
\toprule
container ($10^6$ insertions, then $10^6$ lookups) & random 64-bit keys & keys below $10^6$ \\
\midrule
\texttt{map<ll, int>} & 1.15\,s & 0.95\,s \\
\texttt{unordered\_map<ll, int>} (identity hash, attackable) & 0.29\,s & 0.12\,s \\
\texttt{unordered\_map<ll, int, custom\_hash>} & 0.37\,s & 0.22\,s \\
\quad the same after \texttt{reserve(1000000)} & 0.27\,s & 0.18\,s \\
\texttt{gp\_hash\_table<ll, int, custom\_hash>} & 0.08\,s & 0.08\,s \\
sort the keys once, then \texttt{lower\_bound} (offline) & 0.27\,s & 0.24\,s \\
\texttt{vector<int>} indexed by the key & --- & 0.003\,s \\
\bottomrule
\end{tabular}
\caption{Measured costs (g++ -O2, best of three runs). Every insertion is \texttt{m[key]++}; half of the lookups hit.}
\label{tab:dicts}
\end{table}

Three conclusions follow. A tree costs four to ten times as much as a hash table, so use \cc{map} only when the order matters. The
policy-based \cc{gp_hash_table}, an open-addressing table, is three to four times faster than \cc{unordered_map} with the same hash;
it needs a good hash even more, since its probing breaks down on structured keys, so always pair it with \cc{custom_hash}. And when the
keys are small integers, an array beats every hash table by more than an order of magnitude.

\begin{code}
#pragma push_macro("int")  // as in cposet: PBDS breaks under #define int
#undef int
#include <ext/pb_ds/assoc_container.hpp>
#pragma pop_macro("int")
__gnu_pbds::gp_hash_table<ll, int, custom_hash> cnt;  // custom_hash: cphash
\end{code}

\usetemplate{Anti-hack hash}{Ctrl+Alt+T, Shift+H}{cphash}

\subsection{Hashing sets and multisets}\label{sec:xorhash}

Hashing can also compare whole collections. Give every possible value $v$ an independent random 64-bit label $h(v)$, and represent a
set by the XOR of its elements' labels. Adding or removing an element is one XOR, and two different sets get the same hash with
probability $2^{-64}$: their symmetric difference is a non-empty set of independent uniform labels, whose XOR is again uniform. The
XOR of a \emph{multiset} sees only the parity of every value's count, which is exactly what ``does every value occur an even number
of times?'' needs; to compare multisets themselves, add the labels instead, modulo $2^{64}$.

The power of the idea is that the hash of a range of an array is a difference of two prefix hashes, so questions about the
\emph{contents} of subarrays become comparisons of numbers.

\begin{figure}[htbp]
\centering
\begin{tikzpicture}
  \tikzset{pa/.style={vtxs, cool}, pb/.style={vtxs, warn}, pc/.style={vtxs, dim}}
  \node[blbl, anchor=east] at (-0.45, 0) {$a$};
  \arrs{A}{0}{0}{1/n, 2/n, 2/n, 1/n, 3/n, 3/n, 1/n, 1/n}
  \spanabove{A-2}{A-3}{$p_1 = p_3$}
  \spanabove{A-4}{A-7}{$p_3 = p_7$}
  \node[blbl, anchor=east] at (-0.45, -1.05) {prefix};
  \foreach \cls/\pst [count=\kk from 0] in {A/pa, B/pb, C/pc, B/pb, A/pa, D/pc, A/pa, B/pb, A/pa} {
    \node[\pst] at (0.7*\kk, -1.05) {\cls};
    \node[lbl] at (0.7*\kk, -1.6) {$p_{\kk}$};
  }
  \node[notel, anchor=west, text width=7.0cm] at (6.5, -0.55) {Class A: no value occurs an odd number of times, so $p = 0$ (4
    prefixes, 6 pairs). Class B: only 1 is odd, $p = h(1)$ (3 prefixes, 3 pairs). C: 1 and 2 odd, $p = h(1) \oplus h(2)$; D: only
    3 odd, $p = h(3)$. Every pair of equal prefixes bounds an even subarray: $6 + 3 = 9$ of them.};
\end{tikzpicture}
\caption{Prefix XOR hashes of $a = 1, 2, 2, 1, 3, 3, 1, 1$. Prefix $p_k$ sits at the boundary after the first $k$ elements; prefixes
of the same class have equal hashes.}
\label{fig:xorhash}
\end{figure}

\begin{problem}{Even subarrays}
Count the subarrays of $a_1, \dots, a_n$ in which every value occurs an even number of times ($n \le 2 \cdot 10^5$, values up to
$10^9$).
\solution Let $p_i$ be the XOR of the labels of $a_1, \dots, a_i$, with $p_0 = 0$. The subarray $a_{l+1}, \dots, a_r$ qualifies
exactly when every value occurs an even number of times in it, that is, barring a false match, when $p_l = p_r$. So the answer is the
number of pairs of equal prefix hashes (Figure~\ref{fig:xorhash}): sort the $n + 1$ prefixes and add $\binom{c}{2}$ for every run of
$c$ equal values.
\end{problem}

\begin{code}
void solve() {  // RandHash (cpxorhash) is pasted above solve()
    int n;
    cin >> n;
    RandHash h;
    vector<uint64_t> px(n + 1, 0);  // px[i]: the XOR of the labels of a[0..i)
    rep(i, n) {
        ll x;
        cin >> x;
        px[i + 1] = px[i] ^ h(x);  // a value seen twice cancels out
    }
    sort(all(px));  // equal prefixes now sit next to each other
    ll ans = 0;
    for (int i = 0, j = 0; i <= n; i = j) {
        while (j <= n && px[j] == px[i]) j++;
        ans += (j - i) * (j - i - 1) / 2;  // pairs among j - i equal prefixes
    }
    print(ans);
}
\end{code}

\usetemplate{Random hashing of sets and multisets (XOR / sum hashing)}{Ctrl+Alt+U, X}{cpxorhash}

\begin{tip}[How likely is a false match?]
A program that compares $m$ pairs of hashes errs with probability at most $m / 2^{64}$. When it looks for \emph{any} equal pair among
$N$ hashes, as above, about $N^2 / 2$ pairs are compared implicitly, and the bound becomes $N^2 / 2^{65}$: about $10^{-9}$ for $N = 2
\cdot 10^5$. With 32-bit labels the expected number of false matches would be $N^2 / 2^{33} \approx 4.7$, so the program would almost
surely fail. Sums of labels are slightly weaker: if every difference between the two multisets' counts is divisible by $2^t$, the
bound is $2^{t - 64}$, which is still below $2^{-46}$ when the counts are below $2^{18}$.
\end{tip}

The same labels answer many other questions in $\Oh{1}$ after a linear precomputation. Is $a[l, r)$ a permutation of $1, \dots, r - l$?
Compare its sum of labels with the precomputed sum of the labels of $1, \dots, r - l$. Are two subarrays rearrangements of each other?
Compare their sums of labels. Do $u$ and $v$ have the same set of neighbours? Compare the XORs of their neighbours' labels. Each
question costs one comparison instead of a sort.

\begin{summary}[Chapter summary]
\begin{itemize}
  \item Sorted containers answer neighbour and extremum questions in $\Oh{\log n}$; pairing two of them, like positions and gaps,
        handles events that split or merge. Erase one copy from a multiset through \cc{find}, and break comparator ties with an id.
  \item The policy-based tree adds $k$-th element and rank queries; a Fenwick tree with descent does the same when the values are
        known in advance. Store pairs for multisets, never \cc{less_equal}.
  \item Treaps keyed by position implement sequences with insertions, deletions, reversals and block moves through split and merge;
        a node's ancestors are the highest nodes of the ranges around it, so the expected depth is at most $2 \ln n$.
  \item Copying the nodes on the walked paths makes a treap persistent; merging by subtree sizes keeps it balanced even when versions
        share nodes. Budget about 46 new nodes per update.
  \item CDQ divide and conquer reduces offline three-dimensional problems, including ones over time, to sweeps with a Fenwick tree,
        in $\Oh{n \log^2 n}$.
  \item A wavelet matrix answers $k$-th smallest and counting queries on ranges online in $\Oh{\log n}$ with one bit per element and
        level.
  \item Hash tables need a randomised hash, \cc{gp_hash_table} is the fastest of them, and arrays beat them all; random 64-bit labels
        hash whole sets (XOR) and multisets (sums).
\end{itemize}
\end{summary}

\begin{practice}
\ppgroup{Sets and order statistics}
\begin{pplist}
  \item Concert Tickets
  \item Towers
  \item Room Allocation
  \item Traffic Lights
  \item Josephus Problem I
  \item Josephus Problem II
  \item Sliding Window Median
  \item Sliding Window Cost
  \item Nested Ranges Count
  \item Salary Queries
  \item List Removals
\end{pplist}
\ppgroup{Treaps, values in ranges and hashing}
\begin{pplist}
  \item Substring Reversals
  \item Reversals and Sums
  \item Cut and Paste
  \item Distinct Values Queries
  \item Subarray Sums II
  \item Sum of Four Values
\end{pplist}
\end{practice}
